Let \(\widetilde{\mathcal B}_Z\) be the maximal block-separable matrix domain on which the ordinary inverse in Zhao's Equation (8) is defined, and let \(q_Z(B)=[B]\).  Then \(q_Z:\widetilde{\mathcal B}_Z\to\mathcal B_Z\) is surjective, and
\[
q_Z(B)=q_Z(C)
\quad\Longleftrightarrow\quad
C=B\operatorname{diag}(t_1,t_0)
\quad\text{for some }t_1t_0\ne0.
\]
For every such diagonal \(T\), every \(c\in\mathcal C(F_1,F_0)\), and every \(B\in\widetilde{\mathcal B}_Z\),
\[
g_{q_Z(BT)}(c)=g_{q_Z(B)}(c),\qquad
\ell_{q_Z(BT)}(c)=\ell_{q_Z(B)}(c),\qquad
w(q_Z(BT))=w(q_Z(B)).
\]
Consequently, for every nonempty \(K\subseteq\mathcal B_Z\), the infima and suprema of each displayed pullback criterion over \(q_Z^{-1}(K)\) equal those over \(K\), and every attained argmin or argmax correspondence on the matrix domain is exactly the inverse image under \(q_Z\) of its quotient correspondence.  This is an exact quotient transfer for the ordinary-inverse domain implicit in Zhao's formula; it neither defines nor attributes a singular Moore--Penrose domain to Zhao.